% ====================================================================
%  ADCAIJ — Advances in Distributed Computing and Artificial
%  Intelligence Journal (Universidad de Salamanca / BISITE)
% --------------------------------------------------------------------
%  Template note: ADCAIJ is open access, CC-BY, and uses its own
%  LaTeX class `adcaij.cls` available from
%  https://revistas.usal.es/index.php/2255-2863
%  For submission, replace the \documentclass{article} with
%  \documentclass{adcaij}. ADCAIJ accepts manuscripts in two
%  formats (Word and LaTeX) and has no strict word limit; typical
%  papers are 10-16 pages including figures and references.
%  ADCAIJ has a distributed-systems and AI-architecture emphasis, so
%  the framing of this paper foregrounds the training-pipeline /
%  toolkit aspects more than the IJBM version.
% ====================================================================
%% STATUS 2026-07-03: Experimental design superseded by the project audit
%% (research_logs/2026-07-03-project-audit-sota-roadmap.md). This venue awaits
%% the P6 on-device/edge-distillation study. Benchmark v0 results now
%% exist (see research_logs/2026-07-03-*) and are quoted below as preliminary
%% context; final tables require the SLCeleb v1 benchmark and the P6
%% distillation/edge runs (KD into ReDimNet-B0/CAM++, INT8 ONNX export).
% ====================================================================
\documentclass[11pt,a4paper]{article}

% Engine guard: fontspec under XeLaTeX/LuaLaTeX (as in the 2026-05-23
% builds); T1/Latin Modern fallback so the draft also compiles with
% pdfLaTeX on machines without xetex/luaotfload.
\usepackage{iftex}
\iftutex
  \usepackage{fontspec}
  \setmainfont{DejaVu Serif}
  \setsansfont{DejaVu Sans}
  \setmonofont{DejaVu Sans Mono}
\else
  \usepackage[T1]{fontenc}
  \usepackage{lmodern}
\fi

\usepackage[a4paper,margin=1in]{geometry}
\usepackage[hidelinks]{hyperref}
\usepackage{booktabs,tabularx,amsmath,microtype,enumitem,parskip}
\usepackage{authblk}

\title{A modular, multi-GPU training pipeline for low-resource
cross-lingual speaker verification: design, ablation, and reference
deployment on Sinhala--Tamil bilingual audio}

\author[1]{[First Author]}
\author[1]{[Second Author]}
\affil[1]{[Department], [University], Sri Lanka. \\
\texttt{[email@institution]}}

\date{}

\begin{document}
\maketitle

\begin{abstract}
\noindent We describe the design, implementation, and empirical
validation of an open-source training pipeline for cross-lingual
speaker verification under low-resource constraints. The pipeline
extends the upstream \texttt{voxceleb\_trainer} codebase with eleven
modular features --- each selectable via a single configuration key
--- spanning multilingual self-supervised front-ends, cross-lingual
fine-tuning with layer-wise learning-rate decay, auxiliary
language-identification heads, score normalisation, and PLDA
scoring. We deploy the pipeline across a heterogeneous cluster of
NVIDIA T4, A10, and A40 GPUs, with distributed data-parallel training
on the T4 pair and single-GPU runs for ablation passes. As a
reference deployment, we evaluate on a two-tier Sinhala--Tamil
benchmark: v0, 527 speakers of read speech from OpenSLR SLR52 and
SLR65 (478 Sinhala, 49 Tamil; 12{,}000 trials per language), and v1,
built on SLCeleb, our 280-speaker in-the-wild corpus (pending data
placement), with single-feature ablations isolating contributions. The toolkit, master experimental protocol,
corpus-preparation script, and aggregation tools are released to
enable reproducible low-resource SV research.
\end{abstract}

\noindent \textbf{Keywords:} speaker verification; cross-lingual
transfer; distributed training; multi-GPU; open-source toolkit;
low-resource languages; ablation study.

\section{Introduction}
\label{sec:intro}

Modern speaker verification (SV) systems are increasingly defined as
much by the engineering of the training pipeline as by the choice of
architecture or loss function. A typical research setup combines
several independently-developed components: a front-end (mel /
SincConv / self-supervised), a backbone (TDNN / ResNet / transformer),
a margin-based loss (AM- / AAM-Softmax), an evaluation backend
(cosine / PLDA), score normalisation, and one or more auxiliary
training objectives. Yet most published systems combine these
components in a single end-to-end commitment, without ablating
contributions, and without releasing the configuration scaffolding
that would let a downstream practitioner enable or disable each lever.

This paper presents the inverse: a training pipeline that decomposes
the SV system into eleven \emph{independently togglable} features
(FEATURE-001 through FEATURE-011 in the codebase), each guarded by a
single YAML or CLI key, each accompanied by a design document, and
each contributing to a unified single-feature ablation. We describe
the architectural decisions that make this modularity work in
practice (Section~\ref{sec:design}), the heterogeneous multi-GPU
deployment we use to run the protocol (Section~\ref{sec:deployment}),
and the empirical validation on a low-resource Sinhala--Tamil corpus
(Section~\ref{sec:experiments}). Contributions:

\begin{enumerate}
    \item A modular SV training pipeline with eleven independently
    togglable features, each composable with all others and each
    defaulting to off so that pre-existing experiment scripts remain
    behaviour-identical.
    \item A heterogeneous multi-GPU deployment strategy that maps
    each protocol stage to the appropriate device class (DDP on
    consumer GPUs, single-GPU for ablations, large-VRAM for
    self-supervised backbones).
    \item A reproducible two-tier Sinhala--Tamil benchmark (v0: 527
    speakers from OpenSLR SLR52/SLR65 read speech; v1: SLCeleb, our
    280-speaker in-the-wild corpus, pending) as the pipeline's
    reference deployment, with corpus preparation and results
    aggregation fully automated.
\end{enumerate}

\section{Related work}
\label{sec:related}

\subsection{Speaker verification architectures}

ECAPA-TDNN \cite{desplanques2020ecapa} is the de-facto reference
architecture for non-self-supervised SV; multilingual self-supervised
encoders (WavLM \cite{chen2022wavlm}, wav2vec~2 \cite{baevski2020wav2vec2})
provide strong out-of-the-box features for low-resource languages.
Margin-based losses (AAM-Softmax \cite{deng2019arcface}) are standard
for embedding training.

\subsection{Cross-lingual transfer and SV-specific eval backends}

Cross-lingual fine-tuning from English-pretrained checkpoints is the
established remedy for low-resource SV. Eval-time backends such as
PLDA \cite{sizov2014unifying} and AS-Norm \cite{matejka2017analysis}
add 5--20\% relative MinDCF reduction at zero training cost.
Adversarial language/channel invariance \cite{ganin2015dann} is
infrastructure-ready in our pipeline but evaluated only as future
work in this paper.

\subsection{Modular training pipelines in the broader AI ecosystem}

Modular training frameworks --- HuggingFace \texttt{transformers},
SpeechBrain, NeMo, fairseq --- have shaped the way modern ML research
is conducted. Our pipeline shares their feature-toggle philosophy but
specialises it for SV: each feature is paired with a design document
articulating the hypothesis, expected effect size, and risk; each is
ablated explicitly in a single command sweep.

\section{System design}
\label{sec:design}

\subsection{Feature toggle as the primary abstraction}

Every shipped feature obeys a strict invariant: \emph{the default-off
configuration is byte-identical to the pre-feature pipeline}. This
constraint --- enforced by design review for every merge --- means a
practitioner can adopt the pipeline incrementally, enabling one feature
at a time, with confidence that earlier experiments remain valid.

Concretely:
\begin{itemize}[leftmargin=*]
    \item Every feature is gated by a single boolean (e.g.\ \texttt{plda:
    true}); the YAML and CLI representations are isomorphic.
    \item Default-off paths execute the original code path with no
    additional branching cost.
    \item Each feature ships with a smoke test verifying behaviour
    parity at the default setting.
\end{itemize}

\subsection{Composition rules}

The eleven features compose into the following layered architecture
(Figure~\ref{fig:layers}, [INSERT\_FIGURE]):

\begin{description}[leftmargin=*,style=nextline]
    \item[Front-end] Mel-filterbank, SincConv (raw-waveform), or
    self-supervised transformer (WavLM / XLS-R / mHuBERT-147).
    \item[Backbone] ECAPA-TDNN, ResNet-SE34V2, RawNet3, MLPMixerSpeaker
    (raw or mel).
    \item[Pooling] Channel-attentive statistical pooling (ECAPA),
    self-attentive pooling (others).
    \item[Training-time auxiliary heads] Multi-task language-ID,
    DANN-style adversarial language/channel.
    \item[Loss] AAM-Softmax, AM-Softmax, ProtoNet, GE2E, Triplet.
    \item[Optimisation] Selective freezing, layer-wise LR decay,
    distributed data-parallel.
    \item[Eval backends] Cosine, PLDA (two-covariance), AS-Norm
    (composable with both).
    \item[Reporting] Per-language EER / MinDCF, pooled metrics,
    multi-seed aggregation.
\end{description}

\subsection{Configuration management and reproducibility}

A single YAML config encodes every choice in the pipeline. Paths
expand environment variables (\texttt{\$\{SL\_SPV\_DATA\_ROOT\}}) for
machine portability. The deterministic-mode switch
(\texttt{deterministic: true}) sets cuDNN deterministic flag,
disables TF32, and enables \texttt{torch.use\_deterministic\_algorithms}
for paper-grade reproducibility.

\section{Multi-GPU deployment}
\label{sec:deployment}

\subsection{Heterogeneous cluster allocation}
\label{sec:deploy:hetero}

The reference deployment uses three machine types in a small
heterogeneous cluster:

\begin{table}[t]
\centering
\caption{Heterogeneous-cluster GPU allocation for the protocol.}
\label{tab:cluster}
\begin{tabularx}{\linewidth}{lXX}
\toprule
\textbf{Machine} & \textbf{VRAM} & \textbf{Workload} \\
\midrule
A40 (single node)    & 48 GB        & Self-supervised front-end configs (FEATURE-001 WavLM / XLS-R / mHuBERT) \\
2$\times$ T4 (DDP)   & 16 GB $\times$ 2 & ECAPA-TDNN baseline, primary protocol; distributed data-parallel \\
A10 (single node)    & 24 GB        & Single-feature ablations, multi-seed reproducibility runs \\
\bottomrule
\end{tabularx}
\end{table}

Distributed data-parallel on the T4 pair uses NCCL backend; the
trainer reuses the upstream \texttt{voxceleb\_trainer} DDP path
(verified after BUGFIX-001 corrected the original kwarg-passing bug
in \texttt{mp.spawn}).

\subsection{Wall-clock budget}
\label{sec:deploy:budget}

The full protocol (three primary configurations $\times$ three seeds
$+$ five single-feature ablations $\times$ one seed) dispatches to
$\sim$3--5 days of wall-clock time across the three machines.

\section{Experiments}
\label{sec:experiments}

\subsection{Reference deployment: two-tier Sinhala--Tamil benchmark}
\label{sec:exp:corpus}

The reference deployment is a two-tier benchmark. Benchmark v0
comprises 527 speakers of read speech: 478 Sinhala from OpenSLR SLR52
and 49 Tamil from OpenSLR SLR65 (speakers with at least 10
utterances). Benchmark v1 is built on SLCeleb, our 280-speaker
in-the-wild Sinhala--Tamil corpus (multi-genre, multi-session), and is
pending data placement. Train/test split is 80/20 within each speaker;
trial pairs (v0: 3{,}000 target + 9{,}000 impostor pairs per language)
are generated deterministically from a single random seed
(\texttt{tools/sl\_dataprep.py}). Per-language, cross-lingual, and
pooled trial lists are evaluated independently; the cross-lingual list
is empty at v0 (no bilingual speakers in OpenSLR) and is populated at
v1 by SLCeleb.

\subsection{Configurations}
\label{sec:exp:configs}

\begin{itemize}[leftmargin=*]
    \item P0: ECAPA-TDNN from scratch.
    \item P1: P0 + cross-lingual fine-tune from VoxCeleb1 ECAPA
    checkpoint, front-end frozen, layer-wise LR decay $\gamma=0.9$.
    \item P2: P1 + multi-task language-ID head ($\lambda=0.3$) +
    AS-Norm (top-$K=300$, 500-speaker cohort) + PLDA ($d=200$) +
    per-language evaluation.
\end{itemize}

\subsection{Preliminary results (benchmark v0)}
\label{sec:exp:prelim}

%% Numbers below are from the 2026-07-03 benchmark v0 runs
%% (research_logs/2026-07-03-{zeroshot-baseline-table,p2-backend-adaptation,
%% p3-peft-finetuning}.md). The P6 distillation/edge results this paper
%% targets are pending.

\begin{table}[t]
\centering
\caption{Preliminary EER (\%) on benchmark v0 (527 read-speech
speakers; 12{,}000 trials per language). Fine-tuned rows are
$\mathrm{mean} \pm \mathrm{std}$ over three seeds. Preliminary,
benchmark v0 (read speech, closed-set for the fine-tuned rows);
distillation and on-device results are pending.}
\label{tab:prelim}
\begin{tabular}{lccc}
\toprule
\textbf{System} & \textbf{Params} & \textbf{Si} & \textbf{Ta} \\
\midrule
Zero-shot SpeechBrain ECAPA-TDNN        & $\sim$20M & 4.78 & 3.21 \\
Zero-shot ReDimNet-B1                   & 2.2M      & 3.57 & 1.67 \\
Zero-shot ReDimNet-B6                   & 15M       & 2.71 & 1.48 \\
ReDimNet-B6 + AS-Norm (training-free)   & 15M       & 2.07 & 0.90 \\
WavLM-base+ LoRA fine-tune (3 seeds)    & 95M (1.1\% trained) & 3.73 $\pm$ 0.49 & 3.11 $\pm$ 0.61 \\
WavLM-base+ full fine-tune (3 seeds)    & 95M       & \textbf{1.69 $\pm$ 0.06} & \textbf{1.41 $\pm$ 0.10} \\
\bottomrule
\end{tabular}
\end{table}

Table~\ref{tab:prelim} gives preliminary benchmark-v0 context for the
edge study: the 2.2M-parameter ReDimNet-B1 already outperforms the
$\sim$20M ECAPA-TDNN zero-shot, AS-Norm is a training-free lever worth
carrying on-device (Si 2.71 $\to$ 2.07), and the fine-tuned
WavLM-base+ (full fine-tune) is the current distillation-teacher
candidate. Per-language score calibration proved mandatory
(cross-language calibrator swap inflates Cllr up to $5.7\times$) ---
an on-device deployment must ship per-language calibrators. The
distillation, quantisation, and runtime results that are this paper's
subject remain pending (P6).

\subsection{Headline results}

\begin{table}[t]
\centering
\caption{Equal-error rate (\%) on the SL-SPV evaluation set,
$\mathrm{mean} \pm \mathrm{std}$ over three random seeds.}
\label{tab:headline}
\begin{tabular}{lcccc}
\toprule
\textbf{Config} & \textbf{Pooled} & \textbf{Si} & \textbf{Ta} & \textbf{Cs} \\
\midrule
P0 from scratch   & [X.X $\pm$ Y.Y] & [X.X $\pm$ Y.Y] & [X.X $\pm$ Y.Y] & [X.X $\pm$ Y.Y] \\
P1 + fine-tune    & [X.X $\pm$ Y.Y] & [X.X $\pm$ Y.Y] & [X.X $\pm$ Y.Y] & [X.X $\pm$ Y.Y] \\
P2 full stack     & \textbf{[X.X $\pm$ Y.Y]} & \textbf{[X.X $\pm$ Y.Y]}
                  & \textbf{[X.X $\pm$ Y.Y]} & \textbf{[X.X $\pm$ Y.Y]} \\
\bottomrule
\end{tabular}
\end{table}

\subsection{Single-feature ablations}

\begin{table}[t]
\centering
\caption{Single-feature ablations from the full stack (P2), seed 42.}
\label{tab:ablation}
\begin{tabular}{lcc}
\toprule
\textbf{Ablation} & \textbf{Pooled EER \%} & $\boldsymbol{\Delta}$ \\
\midrule
P2 (all features)   & [X.X] & 0 \\
$-$ fine-tune       & [X.X] & $+$[Y.Y] \\
$-$ LLRD            & [X.X] & $+$[Y.Y] \\
$-$ lang-aux        & [X.X] & $+$[Y.Y] \\
$-$ AS-Norm         & [X.X] & $+$[Y.Y] \\
$-$ PLDA            & [X.X] & $+$[Y.Y] \\
\bottomrule
\end{tabular}
\end{table}

\section{Discussion}
\label{sec:discussion}

\subsection{Modularity in practice}

[NOTE: discuss whether the strict default-off invariant held in
practice across the protocol; identify any place where a feature's
behaviour leaked into the default path; reflect on the trade-off
between modularity and absolute SOTA performance.]

\subsection{Heterogeneous GPU utilisation efficiency}

[NOTE: report wall-clock utilisation and idle-time fraction per
machine; recommend allocation revisions if a particular workload
turned out to be limited by something other than VRAM/compute.]

\subsection{Limitations}

[NOTE: benchmark v0 is read speech, near single-session, and
closed-set for the fine-tuned systems (optimistic absolute EERs;
benchmark v1 / SLCeleb provides the open-set in-the-wild condition);
single-deployment generalisation; absence of self-supervised continued
pretraining (FEATURE-009 is infrastructure-ready but not exercised
here).]

\section{Conclusion}
\label{sec:conclusion}

We described and validated a modular, multi-GPU training pipeline for
low-resource cross-lingual speaker verification. The pipeline's eleven
independently-togglable features compose cleanly under default-off
invariants and ablate cleanly under single-feature negations. The
reference deployment on Sinhala--Tamil bilingual audio demonstrates the
pipeline's intended use; the open-source release positions the
pipeline as the substrate for downstream low-resource SV research.

\section*{Acknowledgements}
[Funding, computational resources, advisor acknowledgements.]

\bibliographystyle{plain}
\bibliography{../references}

\end{document}
